% ============================================================================
% PRACTITIONER'S GUIDE: Using Semantic Transfer in Production
% ============================================================================
% Actionable guidance for deploying banditGPT with semantic transfer
% Based on regime-dependent n_eff sensitivity analysis
% ============================================================================

\subsection{Practitioner's Guide: Deploying Semantic Transfer}
\label{sec:practitioners_guide}

Our sensitivity analysis reveals an important principle for production deployment: \textbf{system robustness comes from Corralling's adaptive expert selection, not from hyperparameter optimization}. This section provides actionable guidance for practitioners deploying banditGPT with semantic transfer.

% ============================================================================
% Quick Reference Box
% ============================================================================

\begin{tcolorbox}[
    colback=blue!5!white,
    colframe=blue!75!black,
    title=\textbf{Quick Reference: Semantic Transfer Deployment}
]

\textbf{Default Configuration (Recommended):}
\begin{itemize}[leftmargin=*, itemsep=1pt]
    \item Set \texttt{n\_effective = 5.0} (mid-range value)
    \item Enable \texttt{use\_corralling = True} (adaptive expert selection)
    \item Let Corralling choose when to use semantic transfer
\end{itemize}

\textbf{Expected Behavior:}
\begin{itemize}[leftmargin=*, itemsep=1pt]
    \item \textbf{33\% of traffic}: Corralling uses semantic transfer (warmup expert)
    \item \textbf{67\% of traffic}: Corralling uses cold-start exploration (tabula rasa expert)
    \item \textbf{Overall impact}: $\sim$1.5\% performance improvement from \texttt{n\_effective} tuning
\end{itemize}

\textbf{Don't waste time on:}
\begin{itemize}[leftmargin=*, itemsep=1pt]
    \item Grid-searching \texttt{n\_effective} (0.33 $\times$ 4.6\% = 1.5\% max benefit)
    \item Forcing semantic transfer (may underperform cold start by 3.9\%)
    \item Manually validating priors (Corralling does this automatically)
\end{itemize}

\end{tcolorbox}

% ============================================================================
% Decision Tree for Practitioners
% ============================================================================

\subsubsection{Decision Tree: Should You Use Semantic Transfer?}

\textbf{Question 1: Is your new model semantically similar to an existing model?}

\begin{itemize}[leftmargin=*, itemsep=2pt]
    \item \textbf{Yes} (e.g., GPT-4-Turbo $\to$ GPT-5.1, Mixtral-8x7B $\to$ Mixtral-8x22B):
    \begin{itemize}
        \item Proceed to Question 2
        \item Check: cosine similarity of embeddings $>$ 0.7
    \end{itemize}
    
    \item \textbf{No} (e.g., text-only pool $\to$ first multimodal model):
    \begin{itemize}
        \item Use cold start (\texttt{n\_effective = 0.5}, identity prior)
        \item Semantic transfer not applicable
    \end{itemize}
\end{itemize}

\textbf{Question 2: Do you trust your warmup priors to match production traffic?}

\begin{itemize}[leftmargin=*, itemsep=2pt]
    \item \textbf{Unsure / Don't know}:
    \begin{itemize}
        \item \textbf{Recommended}: Enable Corralling (\texttt{use\_corralling = True})
        \item Corralling will automatically detect prior mismatch
        \item Expected: 33-67\% split between warmup and tabula rasa experts
        \item Monitor expert selection frequencies (see Section~\ref{sec:monitoring})
    \end{itemize}
    
    \item \textbf{Yes, confident priors match}:
    \begin{itemize}
        \item Use semantic transfer (\texttt{n\_effective = 1.0} for weak priors)
        \item Optionally disable Corralling for faster routing
        \item Warning: May underperform by 3.9\% if priors mismatch (Table~\ref{tab:ablation_no_corralling})
    \end{itemize}
    
    \item \textbf{No, priors likely mismatch}:
    \begin{itemize}
        \item Either: Enable Corralling (will auto-switch to tabula rasa)
        \item Or: Use cold start directly (\texttt{n\_effective = 0.5})
    \end{itemize}
\end{itemize}

% ============================================================================
% Monitoring Guidance
% ============================================================================

\subsubsection{What to Monitor in Production}
\label{sec:monitoring}

Don't obsess over \texttt{n\_effective}. Instead, monitor \textbf{expert selection behavior}:

\paragraph{Key Metrics:}
\begin{enumerate}[leftmargin=*, itemsep=2pt]
    \item \textbf{Expert Weight Distribution}:
    \begin{itemize}
        \item Track: \texttt{corralling\_router.weights[0]} (warmup expert weight)
        \item Healthy range: 30-70\% (adaptive switching active)
        \item Red flag: 100\% either expert for extended periods (>1000 steps)
    \end{itemize}
    
    \item \textbf{Regime Frequency}:
    \begin{itemize}
        \item Expected: $\sim$33\% warmup-dominant, $\sim$67\% tabula rasa-dominant
        \item Measure: Fraction of time warmup weight $>$ 50\%
        \item Diagnostic: If 100\% warmup, investigate prior quality (may be overconfident)
    \end{itemize}
    
    \item \textbf{Performance Stratified by Regime}:
    \begin{itemize}
        \item Report: Mean reward when warmup expert is active vs tabula rasa active
        \item Useful for: Understanding when priors help vs hurt
        \item Action: If warmup regime underperforms tabula rasa, retrain priors
    \end{itemize}
\end{enumerate}

\paragraph{Red Flags (When to Intervene):}
\begin{itemize}[leftmargin=*, itemsep=2pt]
    \item \textbf{100\% warmup weight + low performance}: Priors are overconfident, forcing underexploration
        \begin{itemize}
            \item Fix: Decrease \texttt{n\_effective} or increase \texttt{corralling\_gamma}
        \end{itemize}
    
    \item \textbf{100\% tabula rasa weight immediately}: Priors are catastrophically wrong
        \begin{itemize}
            \item Fix: Retrain warmup priors on more representative data
        \end{itemize}
    
    \item \textbf{Frequent switching (weight oscillates)}: Priors are marginally useful
        \begin{itemize}
            \item Fix: Increase \texttt{corralling\_learning\_rate} for faster commitment
        \end{itemize}
\end{itemize}

% ============================================================================
% Common Pitfalls
% ============================================================================

\subsubsection{Common Pitfalls to Avoid}

\paragraph{Pitfall 1: Forcing semantic transfer when priors mismatch.}
\textit{Symptom}: Performance degrades after new model release despite semantic similarity.

\textit{Root cause}: Warmup priors encode ``expensive-bias'' from historical data, but production traffic consists of simple prompts where expensive models offer no quality advantage.

\textit{Solution}: Enable Corralling. It will detect the mismatch and switch to cold-start exploration (67\% of cases in our data).

\paragraph{Pitfall 2: Obsessing over \texttt{n\_effective} tuning.}
\textit{Symptom}: Spending days grid-searching \texttt{n\_effective $\in$ \{0.5, 1.0, 2.0, 5.0, 10.0, 20.0\}}.

\textit{Root cause}: Misunderstanding that robustness comes from Corralling's adaptive behavior, not parameter insensitivity.

\textit{Solution}: Use default \texttt{n\_effective = 5.0} and trust Corralling. Expected benefit from optimization: $\sim$1.5\% (0.33 $\times$ 4.6\%).

\paragraph{Pitfall 3: Disabling Corralling for ``simplicity.''}
\textit{Symptom}: Router uses pure semantic transfer (single expert), performance is brittle to prior quality.

\textit{Root cause}: Without Corralling, system cannot adapt when priors fail. Can underperform cold start by 3.9\% (Table~\ref{tab:ablation_no_corralling}).

\textit{Solution}: Keep Corralling enabled. Overhead is minimal (2-3 experts, importance-weighted loss computation).

% ============================================================================
% FAQ
% ============================================================================

\subsubsection{Frequently Asked Questions}

\paragraph{Q: Should I tune \texttt{n\_effective} differently per model?}
\textbf{A}: No. Use a global default (\texttt{n\_effective = 5.0}). Model-specific tuning offers marginal gains ($<$2\%) and complicates configuration management.

\paragraph{Q: When should I retrain warmup priors?}
\textbf{A}: Monitor expert selection. If Corralling consistently uses tabula rasa expert (>80\% of time), priors are outdated. Retrain on recent production data (last 30 days recommended).

\paragraph{Q: Can I disable semantic transfer for certain models?}
\textbf{A}: Yes. Set \texttt{n\_effective = 0.5} for that specific model (equivalent to cold start). Useful for dissimilar models (e.g., first multimodal model in text-only pool).

\paragraph{Q: What if I don't have enough data for warmup priors?}
\textbf{A}: Use cold start (\texttt{priors = None}, \texttt{n\_effective = 0.5}). Alternatively, transfer from a semantically similar public model (e.g., use GPT-3.5 priors for Llama-2).

\paragraph{Q: How do I know if 33\%/67\% split applies to my data?}
\textbf{A}: It's data-dependent. Our 67\% tabula rasa frequency reflects 71.5\% ties in LMSYS test set (low task variance). Domain-specific applications (e.g., legal, medical) may have higher task variance $\to$ more warmup usage. Monitor your production split.

% ============================================================================
% Design Principles
% ============================================================================

\subsubsection{Design Principles (Summary)}

\begin{enumerate}[leftmargin=*, itemsep=2pt]
    \item \textbf{Trust meta-learning over manual tuning}: Let Corralling decide when to use semantic transfer. It adapts automatically to data-prior match quality.
    
    \item \textbf{Monitor behavior, not hyperparameters}: Track expert selection frequencies, performance stratified by regime, not \texttt{n\_effective} values.
    
    \item \textbf{Selective application beats universal reliance}: Robustness comes from using semantic transfer only when appropriate (33\% of traffic), not forcing it everywhere.
    
    \item \textbf{Weak priors preserve exploration}: When semantic transfer is used, prefer \texttt{n\_effective = 1.0} over strong priors (>10.0) to avoid over-confidence trap.
    
    \item \textbf{Cold start is not a failure state}: When priors mismatch, cold-start exploration (tabula rasa expert) is the \textit{correct} strategy, not a fallback.
\end{enumerate}

% ============================================================================
% Closing Recommendation
% ============================================================================

\paragraph{Recommended Configuration for Production:}

\begin{verbatim}
router = BanditRouter.create(
    model_registry={...},
    priors="path/to/warmup_priors.joblib",   # Trained on 40k battles
    use_corralling=True,                      # Enable adaptive expert selection
    n_effective_default=5.0,                  # Mid-range default
    corralling_learning_rate=0.1,             # Standard meta-learning rate
    corralling_gamma=0.05,                    # Regularization for stability
    alpha=2.0                                 # Exploration parameter (LinUCB)
)
\end{verbatim}

This configuration:
\begin{itemize}[leftmargin=*, itemsep=1pt]
    \item Enables semantic transfer with reasonable confidence
    \item Allows Corralling to abandon transfer when priors fail
    \item Requires no manual hyperparameter tuning
    \item Achieves robust performance across diverse traffic patterns
\end{itemize}

\textbf{Key insight}: Optimization effort should focus on \textit{prior quality} (retraining on representative data) and \textit{monitoring expert selection}, not on tuning \texttt{n\_effective}.

\vspace{0.3cm}
\hrule
\vspace{0.2cm}
\textit{For implementation details, see \texttt{src/bandit\_gpt/router.py}. For reproduction, see \texttt{experiments\_v1/08\_figure/run\_figure8\_analysis.py}.}
